\section{Discussion and Future Work}
\subsection{Interpreting recovery}
This run did not establish the failure needed to answer the intended recovery
question. The reviewer followed the evidence requirement throughout formation,
despite the conflicting reward and the accepting peer. Accurate later review
therefore cannot demonstrate reversal of collusion. The zero correction and
reset contrasts also cannot show that these interventions would be unnecessary
or ineffective after a harmful convention had formed.

The observed distinction was between accurate classification with records and
blanket rejection without them. False approval alone would miss this difference:
it was zero in every condition. Correct approval distinguishes useful review from
rejection required by missing evidence. Neither observed accuracy nor low false
approval certifies independent reasoning or deployment safety.

The contrast with earlier collusion findings~\citep{shi2026emergent} does not identify why induction
was absent here. Model choice, a scripted neutral peer, short interactions,
repeated verification instructions, and the absence of tools or reflection all
differ from broader agent workflows. This experiment does not isolate those
explanations. Appendix~\ref{app:ethics} states the limits and release disclosures.

\subsection{Priorities for further work}
The next experiment should first establish a failure influenced by peer behavior.
Model and infrastructure comparisons should
use fresh histories and fixed sampling plans, with all assigned outcomes
reported. Recovery remains the central question; establishing the failure makes
that question testable.

\textbf{Does correction last?} The current correction rejects prior shortcuts and
reciprocal approval. Compare this specific correction with a general prohibition
and an explicit verification procedure. Match instruction length and authority.

Resume peer interaction on new claims after a one-time correction, repeated
reminders, or scheduled reminder withdrawal. Track false and correct approval
rates over time, and define recurrence in advance. This tests whether recovery
survives renewed interaction.

\textbf{What carries the behavior forward?} Feedback and memory can affect
collusion~\citep{shi2026emergent}. Remove peer messages, observed verdicts, score
feedback, or stored summaries one at a time. Preserve task facts and match context
length.

Reintroduce each component after a reset to test whether failures return. Compare
reset alone with correction alone. Errors without peer messages do not establish
hidden communication: action history or shared feedback may still shape behavior.

\begin{samepage}
\textbf{When does recovery need enforcement?} Compare correction with an
independent checker that requires complete evidence and verifies answers before
approval. Record proposed verdicts separately from enforced outcomes, then remove
the checker. This distinguishes blocking errors from changing agent behavior.\par
\end{samepage}

Measure correct approvals, latency, and checking cost alongside false approvals.
Rejecting every claim is not useful recovery. Because exact checking is easy in
this synthetic task, also test settings where it is costly.

\textbf{Does recovery generalize?} Use more independent histories, a fixed sample
size, and predefined absolute recovery criteria. Vary model family and agent
infrastructure separately, then test adaptive peers, agent-generated work, and
sandboxed shared artifacts.

Preserve legitimate task information when removing communication channels.
These tests would assess transfer while distinguishing peer influence from
independent adaptation.

\section{Conclusion}
The controlled histories did not induce unsupported approval in this run.
Full-evidence review was accurate, restricted-evidence review rejected every
claim, and correction and reset produced no observed added difference. These
findings leave recovery from established collusion unanswered. A successful
test must first demonstrate the harmful peer-dependent behavior and then assess
whether an intervention restores useful verification.
